\documentclass[11pt]{article}

\usepackage{amsmath}
\usepackage{amssymb}
\usepackage{array}
\usepackage{bm}
\usepackage{esvect}
\usepackage{geometry}
\usepackage{makecell}
\usepackage{siunitx}

\geometry{margin=1in}

\DeclareRobustCommand{\myvec}[1]{\vv{\mathbf{#1}}}
\newcommand{\myunitvec}[1]{\bm{\hat{\mathbf{#1}}}}
\newcommand{\myvecsp}[1]{\bm{\mathcal{#1}}}
\newcommand{\mysubvecsp}[1]{\bm{{#1}}}
\newcommand{\myset}[1]{\mathcal{#1}}
\newcommand{\mymat}[1]{#1}

\title{Diagonalization of Matrices}
\author{\texttt{O5-37-2718281828459}}

\raggedright
\begin{document}
    \subsubsection{The way for a Matrix to record a Linear Transformation}
    Under $\mathbb{F}^n$, using the basis $\myset{B} = \{ \myvec{b}_1, \myvec{b}_2, \dots, \myvec{b}_n \}$ \vspace{1ex}

    Let $T = \Big[ \myvec{t}_1\,\myvec{t}_2\,\dots\,\myvec{t}_n \Big]$ be a $m \times n$ matrix and $L$ be a linear transformation function such that
    \[ \mymat{L}(\myvec{x}) = \mymat{T}\myvec{x} \qquad \text{ for any }n \times 1\text{ column vector }\myvec{x} \]
    
    The Matrix $\mymat{T}$ records $\mymat{L}$ by tracking how the vectors being transformed, \textbf{relatively} \\
    Since $\mymat{L}$ is a linear transformation, \\
    \begin{itemize}
        \item the transformation~applied to the vectors, or the motion~of~the~vectors, can be recorded simply just by the state of vectors at the two ends of the transformation, i.e. the original~vectors~$\myvec{u}$ and the transformed~vectors~$\myvec{v}$
        \item the space is evenly-stretch, shrink or rotated, so the transformation is independent to any specific vector, or, the same transformation rule is applied to all vectors
    \end{itemize} 
    Since we actually does not care about any specific vectors, we just care about the relative transformation, we can directly set the $\myvec{u}$s to be a set of vectors = to the basis vectors, \vspace{1ex}

    so they will be always $\{
        \begin{bmatrix} 1 \\ 0 \\ \vdots \\ 0 \end{bmatrix}, 
        \begin{bmatrix} 0 \\ 1 \\ \vdots \\ 0 \end{bmatrix}, 
        \dots, 
        \begin{bmatrix} 0 \\ 0 \\ \vdots \\ 1 \end{bmatrix} 
    \}$ under the basis, and then the $\myvec{v}$s will be the matrix 
    \begin{equation}
        \mymat{T} = \Big[ [\mymat{L}(\myvec{b}_1)]_{\myset{B}} \Big]
    \end{equation}
    
\end{document}